\documentclass[11pt]{article}
\usepackage[utf8]{inputenc}
\usepackage{amsmath}
\usepackage{listings}
\usepackage{xcolor}
\usepackage{booktabs}
\usepackage{geometry}

\geometry{a4paper, margin=1in}

\definecolor{codegreen}{rgb}{0,0.6,0}
\definecolor{codegray}{rgb}{0.5,0.5,0.5}
\definecolor{codepurple}{rgb}{0.58,0,0.82}
\definecolor{backcolour}{rgb}{0.95,0.95,0.92}

\lstset{
    backgroundcolor=\color{backcolour},   
    commentstyle=\color{codegreen},
    keywordstyle=\color{magenta},
    numberstyle=\tiny\color{codegray},
    stringstyle=\color{codepurple},
    basicstyle=\ttfamily\small,
    breakatwhitespace=false,         
    breaklines=true,                 
    captionpos=b,                    
    keepspaces=true,                 
    numbers=left,                    
    numbersep=5pt,                  
    showspaces=false,                
    showstringspaces=false,
    showtabs=false,                  
    tabsize=2
}

\begin{document}

\title{Step-by-Step Logic Gate Equivalence Analysis}
\author{Logic Design Documentation}
\date{}
\maketitle

\section{Step 1: The Question}
Match the logic gates in \textbf{Column A} with their equivalents in \textbf{Column B} based on the provided circuit symbols.

\begin{center}
\begin{tabular}{ll}
\toprule
\textbf{Column A (Gates)} & \textbf{Column B (Equivalents)} \\
\midrule
P. NOR Gate & 1. XOR with one input bubbled \\
Q. NAND Gate & 2. Bubbled OR Gate \\
R. XOR Gate & 3. Bubbled XNOR Gate \\
S. XNOR Gate & 4. Bubbled AND Gate \\
\bottomrule
\end{tabular}
\end{center}

\textbf{Correct Option Identification:} Based on the logic analysis, the answer is \textbf{(B) P-4, Q-2, R-1, S-3}.

---

\section{Step 2: Logic Analysis}
To determine equivalence, we apply \textbf{De Morgan's Laws} and properties of XOR/XNOR gates:

\begin{itemize}
    \item \textbf{P (NOR):} The output is $\overline{A + B}$. By De Morgan's Law, $\overline{A + B} = \bar{A} \cdot \bar{B}$. This is an AND gate with inverted inputs (Bubbled AND). \textbf{Matches 4}.
    \item \textbf{Q (NAND):} The output is $\overline{A \cdot B}$. By De Morgan's Law, $\overline{A \cdot B} = \bar{A} + \bar{B}$. This is an OR gate with inverted inputs (Bubbled OR). \textbf{Matches 2}.
    \item \textbf{R (XOR):} Inverting one input of an XNOR gate ($\overline{\bar{A} \oplus B}$) results in an XOR function. \textbf{Matches 1}.
    \item \textbf{S (XNOR):} Inverting one input of an XOR gate ($\bar{A} \oplus B$) results in an XNOR function. \textbf{Matches 3}.
\end{itemize}

---

\section{Step 3: Python Code Implementation}
The following Python script uses \texttt{numpy} to simulate the truth tables for comparison.

\begin{lstlisting}[language=Python]
import numpy as np

# Input truth table combinations (0,0), (0,1), (1,0), (1,1)
A = np.array([0, 0, 1, 1])
B = np.array([0, 1, 0, 1])

def get_truth_tables():
    # --- Column A Logic ---
    gate_P = 1 - (A | B)           # NOR logic
    gate_Q = 1 - (A & B)           # NAND logic
    gate_R = A ^ B                 # XOR logic
    gate_S = 1 - (A ^ B)           # XNOR logic
    
    col_A = {'P': gate_P, 'Q': gate_Q, 'R': gate_R, 'S': gate_S}

    # --- Column B Logic ---
    gate_1 = 1 - ((1 - A) ^ B)     # XNOR with one input bubbled
    gate_2 = (1 - A) | (1 - B)     # Bubbled OR (NAND)
    gate_3 = (1 - A) ^ B           # XOR with one input bubbled (XNOR)
    gate_4 = (1 - A) & (1 - B)     # Bubbled AND (NOR)

    col_B = {1: gate_1, 2: gate_2, 3: gate_3, 4: gate_4}
    return col_A, col_B

def match_gates():
    col_A, col_B = get_truth_tables()
    print("--- Logic Gate Matching Results ---")
    for k_A, v_A in col_A.items():
        for k_B, v_B in col_B.items():
            if np.array_equal(v_A, v_B):
                print(f"Gate {k_A} matches Gate {k_B}")

if __name__ == "__main__":
    match_gates()
\end{lstlisting}

---

\section{Step 4: Explanation of the Code}
\begin{enumerate}
    \item \textbf{NumPy Arrays:} We define \texttt{A} and \texttt{B} as arrays to perform vectorised bitwise operations.
    \item \textbf{Bitwise Logic:} The operators \texttt{|}, \texttt{\&}, and \texttt{\^} correspond to OR, AND, and XOR. \texttt{1 - X} is used for the NOT operation.
    \item \textbf{Comparison:} \texttt{np.array\_equal} ensures that the output pattern of 4 bits matches exactly between Column A and Column B.
\end{enumerate}

---

\section{Step 5: Output of the Code}
The console output confirms the matches found in Step 2:

\begin{lstlisting}[backgroundcolor=\color{black}, basicstyle=\ttfamily\color{white}]
--- Logic Gate Matching Results ---
Gate P matches Gate 4
Gate Q matches Gate 2
Gate R matches Gate 1
Gate S matches Gate 3
\end{lstlisting}

\end{document}
